% Generated draft: requires mathematical and source review.
\documentclass[11pt]{article}
\usepackage[margin=25mm]{geometry}
\usepackage{fontspec}
\setmainfont{Latin Modern Roman}
\usepackage{amsmath,amssymb,mathtools,unicode-math}
\setmathfont{Latin Modern Math}
\usepackage{enumitem,needspace,xurl,hyperref}
\setlength{\parindent}{0pt}
\setlength{\parskip}{5pt}
\setlength{\emergencystretch}{3em}
\newcommand{\N}{\mathbb{N}}
\newcommand{\Z}{\mathbb{Z}}
\newcommand{\Q}{\mathbb{Q}}
\newcommand{\R}{\mathbb{R}}
\newcommand{\C}{\mathbb{C}}
\newcommand{\F}{\mathbb{F}}
\newcommand{\A}{\mathbb{A}}
\newcommand{\PP}{\mathbb{P}}
\newcommand{\E}{\mathbb{E}}
\newcommand{\eps}{\varepsilon}
\newcommand{\dd}{\,\mathrm{d}}
\begin{document}
\section*{UnsolvedMath ID: 11000315}
\textbf{Source problem code:} catalogue entry
\subsection{Definitions and mathematical statement}
Let \(F_n\) be the rank-\(n\) free group and \(\operatorname{Out}(F_n)=\operatorname{Aut}(F_n)/\operatorname{Inn}(F_n)\). The conjecture is vanishing of stable reduced rational cohomology:
\[\lim_{n\to\infty}\widetilde H^*(\operatorname{Out}(F_n);\Q)=0.\]
In the degreewise stable interpretation this means
\[\forall k\geq1\ \exists N_k\ \forall n\geq N_k,\qquad
H^k(\operatorname{Out}(F_n);\Q)=0.\]
It does not demand one rank threshold working simultaneously in every degree. Ordinary degree-zero cohomology remains \(\Q\); reduced degree zero vanishes. There is no positive-degree exception, in particular none at degree one. Stable maps/ranges are those of the intended homological-stability framework, not an unstated natural inclusion \(\operatorname{Out}(F_n)\hookrightarrow\operatorname{Out}(F_{n+1})\).

The source then proposes studying cohomology with twisted coefficients. A precise version must specify coefficient modules \(V_n\), the base ring and, if a stable question is intended, compatibility as \(n\) varies, and ask for \(H^k(\operatorname{Out}(F_n);V_n)\). The trivial-coefficient conjecture does not assert vanishing for arbitrary such modules or supply a particular family to compute.

The supplied summary attributes the stable rational result to Galatius's stable-homology theorem and marks it solved. Its tentative phrases ``mostly,'' a possible degree-one exception and its informal sphere-space notation do not alter the exact all-positive-degree source assertion. The precise integral infinite-loop-space model is not reconstructed from that imprecise summary.
\subsection{Short English statement}
Stable rational cohomology should vanish in every positive degree, with only ordinary degree zero remaining. The record reports this solved by Galatius, while the separate twisted-coefficient programme needs specified modules.
\subsection{Sources}
\begin{itemize}
\item Local UnsolvedMath registry, permanent ID \(11000315\), exact source code \texttt{catalogue entry}.
\item Farb, \emph{Problems on Mapping Class Groups and Related Topics}, Conjecture 4.9, PDF page 366: \url{https://www.math.uchicago.edu/~farb/papers/mcgbook.pdf}. Registry access 2026-07-29; pdf-font-structure extraction.
\item Galatius, \emph{Stable homology of automorphism groups of free groups}, as supplied by the research summary; exact publication details and the Out/Aut stability comparison are not provided.
\end{itemize}
\end{document}
